Since $Q(u) \ne 0$, $x = \frac{B(x,u)}{Q(u)} u + x_{\perp}$ with $B(x_{\perp}, u) = 0$, which gives $V = \R u \oplus u^{\perp}$. By \claimref{std-algebra.inverse-of-vector}, $u^{-1} = u/Q(u)$. Then $u u u^{-1} = u$, and $u x_{\perp} = -x_{\perp} u$ by \claimref{std-algebra.orthogonal-anticommute}, so $u x_{\perp} u^{-1} = -x_{\perp}$. Hence $-u x u^{-1} = -\frac{B(x,u)}{Q(u)}u + x_{\perp} = x - 2\frac{B(x,u)}{Q(u)} u$. The right-hand side is a linear map fixing $u^{\perp}$ and negating $u$, so it is an involution with determinant $-1$; radical vectors lie in $u^{\perp}$. It is an isometry because $Q(-u x u^{-1}) = (u x u^{-1})^{2} = u x^{2} u^{-1} = Q(x)$ by \claimref{std-algebra.vector-squares}. It depends on $u$ only through the line $\R u$. Finally, with $d = a - b$, $B(a, d) = Q(a) - B(a, b) = \tfrac12 Q(d)$ because $Q(a) = Q(b)$, so $s_d(a) = a - d = b$.
